\documentclass[a4paper,11pt]{article}
%% Compiles both from this folder (pdflatex tp5.tex) and from the root of the project
%% (Overleaf: Menu > Main document = tp5_pacbayes_votes/tp5.tex).
\makeatletter
\IfFileExists{tp5_pacbayes_votes/tp5.tex}{\gdef\tpdir{tp5_pacbayes_votes/}\gdef\tpcommon{common/}}{\gdef\tpdir{}\gdef\tpcommon{../common/}}
\edef\input@path{{\tpdir}{\tpcommon}}
\makeatother
\input{tpsetup}
\renewcommand{\polytitle}{Statistical Supervised Machine Learning -- TP 5}
\graphicspath{{figures/}{tp5_pacbayes_votes/figures/}{../common/logos/}{common/logos/}}

\begin{document}
\sloppy
\sheethead{TP 5 -- PAC-Bayes for majority votes: certify, then learn}{%
Last of five practical sessions (1: protocol, 2: bagging, 3: boosting, 4 and 5: PAC-Bayes). This session follows
Sections~15 to~17 of the lecture notes. Bagging, random forests and boosting all output a majority vote; TP~4 certified
the Gibbs classifier, which draws one voter at random, and we now certify the \emph{vote} itself. We first relate the two
with the first-order bound and see that it ignores how the errors of the voters are correlated; the joint error and the
disagreement of pairs of voters then give the tandem bound, which we study on voters whose correlation we control. We
then certify the random forest of the running example of the notes \emph{without any test set}, with its out-of-bag
examples, and we finally \emph{learn} the weights of the forest by minimizing a bound: we derive the self-bounding
algorithm of Section~17 step by step, implement it, and see why the bound to minimize must take the diversity of the
voters into account. The session lasts \textbf{three hours}; Part~\ref{p5:far} contains the difficult questions.}

\begin{sessionbox}[About this session -- TP 5: PAC-Bayes for majority votes]
\textbf{Goal.} By the end of the session you should be able to compute the Gibbs risk, the joint error and the disagreement
of a vote, prove and compute the first-order and tandem bounds, certify a random forest with its out-of-bag examples,
derive a self-bounding algorithm from a bound, implement it, and explain why minimizing the first-order bound destroys the
diversity of a vote.\par
\textbf{To reread before the session.} TP~2 (random forests, Part~4) and TP~4, and in the lecture notes: Section~15 (from
the Gibbs classifier to the majority vote: Definition~15.1, Propositions~15.2, 15.3, 15.5 and~15.6), Section~16 (learning
with the bounds, out-of-bag examples) and Section~17 (deriving a self-bounding algorithm, Propositions~17.1 and~17.2,
Algorithm~13). The exercises of Sections~15 to~17 of the exercise booklet are a good preparation.\par
\textbf{Timing (3 hours).}\quad
\begin{tabular}[t]{ll}Part~1 & 25 min\\
Part~2 & 45 min\\
Part~3 & 45 min\\
Part~4 & 65 min\\
Part~5 \hard & at home\end{tabular}\par
\textbf{To hand in.} Your notebook with the tables of Parts~1 to~4, the figures of Parts~2 and~4, and your proofs.\par
\textbf{Tools.} \code{numpy}, \code{scipy}, \code{matplotlib}, \code{DecisionTreeClassifier}; the toolbox
\code{common/pbtools.py} provides the bounds of the course (\code{fo\_bound}, \code{tnd\_bound}) and Algorithm~13
(\code{tandem\_fixed\_point}) to check your own functions.
\end{sessionbox}

\medskip
\noindent\textbf{Notation (Section 15).} $\Hcal$ is a finite set of voters $h:\mathcal X\to\{-1,+1\}$ and $Q$ a
distribution on $\Hcal$. For an example $(x,y)$, $W_Q(x,y)=\sum_hQ(h)\ind{h(x)\neq y}$ is the $Q$-weighted proportion of
voters that err (Definition~15.1). We use
\begin{itemize}
\item the \textbf{Gibbs risk} $\RD(\GQ)=\E[W_Q]=\sum_hQ(h)\RD(h)$ and the risk of the \textbf{majority vote}
      $\RD(\BQ)=\P\big(\BQ(X)\neq Y\big)$, where $\BQ(x)=\sign\big(\sum_hQ(h)h(x)\big)$ (ties counted as errors);
\item the \textbf{joint error} $\eD(Q)=\sum_{h,h'}Q(h)Q(h')\,\P\big(h(X)\neq Y,\,h'(X)\neq Y\big)$ and the
      \textbf{disagreement} $\dD(Q)=\sum_{h,h'}Q(h)Q(h')\,\P\big(h(X)\neq h'(X)\big)$ (Definition~15.4).
\end{itemize}
The empirical versions are obtained by replacing $\D$ by a sample. As in TP~4, $\klup(q,\varepsilon)$ is the upper
inverse of the binary kl, and $\delta=0.05$.

\begin{tip}
Store the predictions of the voters on a sample in a matrix \code{V} of shape (voters, examples) and the errors in
\code{E = (V != y)}. Then the empirical risks are \code{E.mean(1)}, the matrix of empirical joint errors is
\code{E @ E.T / n}, the Gibbs risk of $Q$ is \code{Q @ L} and its joint error \code{Q @ Lt @ Q}.
\end{tip}

%% =====================================================================
\tppart{Gibbs classifier, majority vote and the first-order bound (Sections 15.1--15.2)}\label{p5:fo}

\begin{enumerate}
\item Show that $\BQ$ errs on $(x,y)$ if and only if $W_Q(x,y)\geq\frac12$, and deduce from Markov's inequality
      (Lemma~13.1) that $\RD(\BQ)\leq2\RD(\GQ)$ (Proposition~15.2).\\
      Combine it with Seeger's bound (TP~4) to obtain a certificate for the vote.
\item On the toy problem of TP~4 (one sample, $m=100$, uniform prior on the $102$ thresholds), compute for the Gibbs
      posteriors $Q_\lambda$, $\lambda\in\{10,30,100,300,3000\}$: the true Gibbs risk, the true risk of the vote (on a fine
      grid of $[0,1]$, since it can be computed exactly), the oracle bound $2\RD(\GQ)$ and the PAC-Bayesian first-order
      bound.\\
      Comment.
\end{enumerate}

%% =====================================================================
\tppart{Condorcet, joint error, disagreement and the tandem bound (Sections 15.3--15.5)}\label{p5:second}

We consider $n=25$ voters with uniform weights and the following model: voter $i$ errs on an example when
$\sqrt c\,Z+\sqrt{1-c}\,\xi_i<\Phi^{-1}(p)$, where $Z$ (the ``difficulty'' of the example) and $\xi_1,\dots,\xi_n$ are
independent $\mathcal N(0,1)$ and $\Phi$ is the c.d.f.\ of $\mathcal N(0,1)$. Every voter has risk $p$, and $c\in[0,1)$
tunes the correlation between their errors ($c=0$: independent errors). Given $Z$, the errors are independent Bernoulli
variables of parameter $\pi(Z)=\Phi\big((\Phi^{-1}(p)-\sqrt cZ)/\sqrt{1-c}\big)$, so all the exact quantities below can be
computed by a one-dimensional integral over $Z$.

\begin{enumerate}
\item \textbf{Condorcet.} For $c=0$, prove Proposition~15.3: $\RD(\BQ)\leq e^{-2n(\frac12-p)^2}$.\\
      Compute the exact risk of the vote for $n\in\{1,5,25,101,1001\}$, $p\in\{0.3,0.45\}$, $c\in\{0,0.3\}$, compare it
      with this bound when $c=0$, and give its limit when $n\to\infty$.\\
      Conclude.
\item \textbf{Identity.} Show that $\RD(\GQ)=\eD(Q)+\frac12\dD(Q)$ (Proposition~15.5; hint: for two voters,
      $\ind{h\neq y}+\ind{h'\neq y}=2\,\ind{h\neq y,\,h'\neq y}+\ind{h\neq h'}$), and that $\eD(Q)=\E[W_Q^2]$.
\item \textbf{Tandem bound.} Show that $\RD(\BQ)\leq\P(W_Q^2\geq\frac14)\leq4\eD(Q)$ (Proposition~15.6).
\item Simulate the errors of the $25$ voters on $m=20\,000$ examples for $p=0.3$, $c=0.3$.\\
      Compute the Gibbs risk, the joint error and the disagreement (a) from $W_Q$, (b) from the pairs of voters (all the
      $n^2$ ordered pairs, including $h=h'$), and the risk of the vote.\\
      Compare with the exact values.
\item For $p\in\{0.2,0.35,0.45\}$ and $c\in[0,0.98]$, plot the risk of the vote and the two oracle bounds:
      first-order $2\RD(\GQ)$ and tandem $4\eD(Q)$.\\
      When is the tandem bound better than the first-order bound? What does it say about diversity?
\end{enumerate}

%% =====================================================================
\tppart{Certifying a random forest without a test set (Sections 15.5 and 16.2)}\label{p5:oob}

We use the running example of Part~IV of the notes: \code{digits}, task ``is the digit smaller than $5$?'', split with
\code{train\_test\_split(X, y, test\_size=0.3, random\_state=0)} ($1257$ training and $540$ test digits), and a forest of
$100$ trees, each learned with \code{DecisionTreeClassifier(max\_features="sqrt")} on a bootstrap sample of \emph{half}
the size of the training set. The prior $P$ is uniform on the \emph{indices} $1,\dots,100$ of the trees, fixed before
the forest is grown. Each tree $h_t$ is evaluated on its \emph{out-of-bag} (OOB) examples $T_t$, the training examples
that are not in its bootstrap sample.

\begin{enumerate}
\item Write \code{grow\_forest(Xtr, ytr, n\_trees, frac, rng)} returning the trees and the boolean $100\times m$ matrix
      of the OOB examples.\\
      Show that an example is out-of-bag for a given tree with probability $(1-\frac1m)^{m/2}\approx e^{-1/2}\approx61\%$,
      and for two given trees with probability $\approx e^{-1}\approx37\%$. Check it.
\item Compute the OOB error $\hat L(h_t)$ of each tree on $T_t$, and the empirical joint error $\hat e(h_t,h_{t'})$ of each
      pair on $T_t\cap T_{t'}$; let $n_1=\min_t|T_t|$ and $m=\min_{t,t'}|T_t\cap T_{t'}|$.\\
      For the uniform posterior $Q=P$, compute the first-order certificate
      $2\klup\big(\sum_tQ(h_t)\hat L(h_t),\frac{\KL(Q\|P)+\ln(2\sqrt{n_1}/\delta)}{n_1}\big)$ and the tandem certificate
      $4\klup\big(\hat e(Q),\frac{2\KL(Q\|P)+\ln(2\sqrt{m}/\delta)}{m}\big)$ (equations~(15.3) and~(15.4), with the OOB
      replacements of Section~16.2).
\item Compute on the test set the risk of the vote, the Gibbs risk, the joint error and the disagreement, and the kl
      test-set bound of the vote (Theorem~13.5).\\
      Compare the three certificates. Which ones need a test set?
\item Why are the OOB certificates valid, although every training example has been used to grow some of the trees?\\
      Why must $m$ be the size of the \emph{smallest} common OOB set?
\end{enumerate}

%% =====================================================================
\tppart{Learning the weights of the forest: deriving a self-bounding algorithm (Sections 16.3 and 17)}\label{p5:learn}

Since the bounds hold for all posteriors $Q$ simultaneously, we may \emph{choose} the weights of the $100$ trees by
minimizing a bound on the OOB data, and keep the value of the bound as a certificate (self-bounding learning,
Section~16.1). We follow the five steps of Section~17.1.

\begin{enumerate}
\item \textbf{The first-order recipe.} With the bound of Theorem~14.4 applied to the OOB errors, the objective is
      $\sum_tQ(h_t)\hat L(h_t)+\frac{\KL(Q\|P)}{\lambda}$, minimized by the Gibbs posterior
      $Q_\lambda(h_t)\propto P(h_t)e^{-\lambda\hat L(h_t)}$ (Proposition~14.6).\\
      For $\lambda\in\{1,2,4,\dots,4096\}$, compute its first-order certificate, its number of \emph{effective trees}
      $1/\sum_tQ(h_t)^2$ and the test risk of its vote.\\
      Comment.
\item \textbf{The tandem objective.} Applying Theorem~14.4 to pairs of voters gives the linear tandem bound
      (Proposition~17.1). Show that, for a fixed $\lambda$, its terms that depend on $Q$ are
      $F_\lambda(Q)=\hat e(Q)+\frac2\lambda\KL(Q\|P)$.
\item \textbf{The optimality condition.} With a Lagrange multiplier for the constraint $\sum_tQ(h_t)=1$
      (Section~4.3), show that a minimizer $Q^\star$ of $F_\lambda$ with positive weights satisfies
      $Q^\star(h)\propto P(h)\,e^{-\lambda\hat e(h,Q^\star)}$, with $\hat e(h,Q)=\sum_{h'}Q(h')\hat e(h,h')$
      (Proposition~17.2).\\
      Interpret: how is a tree penalized, compared with the Gibbs posterior?
\item \textbf{The algorithm.} Implement the damped fixed-point iterations of Algorithm~13 ($\alpha=0.1$) and check your
      function against \code{tandem\_fixed\_point} of \code{common/pbtools.py}.\\
      For $\lambda\in\{1,2,4,\dots,4096\}$, compute the tandem certificate of the output, its number of effective trees and
      the test risk of its vote.\\
      Plot the certificates and the test risks of the two algorithms against $\lambda$, and the weights of the trees
      against their OOB errors for the selected $\lambda$.
\item Why may $\lambda$ be chosen by minimizing the kl certificate, without any union bound? Is the certificate still valid
      if the iterations are stopped before convergence?
\item Apply the two algorithms (with $\lambda$ selected on the certificate) to \code{spambase}, \code{wdbc} and
      \code{australian} (split $70\%/30\%$, \code{random\_state=0}, same forests).\\
      Report, for the uniform weights and for the two learned posteriors, the certificate, the number of effective trees and
      the test risk of the vote.\\
      Which bound is the right objective for learning a vote?
\end{enumerate}

%% =====================================================================
\tppart{Going further \hard}\label{p5:far}

\begin{enumerate}
\item \textbf{The factor 2 cannot be improved.} Take $n$ voters ($n$ odd) with uniform weights, a set $A$ of inputs of
      probability $\frac{2rn}{n+1}$ on which exactly $\frac{n+1}2$ voters err, and no error outside $A$.\\
      Compute $\RD(\GQ)$, $\eD(Q)$, $\RD(\BQ)$, the ratio $\RD(\BQ)/\RD(\GQ)$ and the tandem bound for $n\in\{3,11,101\}$,
      $r\in\{0.05,0.2\}$.\\
      Conclude.
\item \textbf{The C-bound.} Using the Chebyshev--Cantelli inequality $\P(Z-\E Z\geq a)\leq\frac{\mathrm{Var}Z}{\mathrm{Var}Z+a^2}$
      for the margin $M_Q=1-2W_Q$ (Section~15.6), prove the C-bound $\RD(\BQ)\leq1-\frac{(1-2\RD(\GQ))^2}{1-2\dD(Q)}$ when
      $\RD(\GQ)<\frac12$. Add it to the figure of Part~\ref{p5:second}, question~5.
\item \textbf{Convexity.} When all the voters are evaluated on the same examples $x_1,\dots,x_m$, show that
      $\hat e(Q)=\frac1m\sum_iW_Q(x_i)^2$, and deduce that $F_\lambda$ is convex, so that the fixed point of
      Algorithm~13 is its global minimum.
\item \textbf{The linear certificate.} For the posterior selected in Part~\ref{p5:learn}, compute the linear tandem
      certificate of Proposition~17.1 with the union bound over the $13$ values of $\lambda$ ($\ln\frac{13}\delta$ instead
      of $\ln\frac1\delta$) and compare it with the kl certificate. Why is it looser?
\end{enumerate}

\end{document}
